\chapter[Introducción]{Introducción}
\label{cp:introduccion}
% Extensión estimada: sin límite; la ficha orienta 6--7 páginas (ver fichas/01-introduccion/ y guía §5).
% Fuentes: Plan de trabajo (Introducción y objetivos), Propuesta (secciones "Introducción y Propósito" y
% "Marco de la Problemática"), PLAN_TFG_GAVIOLA_FACUNDO.pdf y la revisión del tutor del 28/09/2026
% (tareas_revision_cap1-2.md, I1--I8). Sin hipótesis (D1). Las preguntas van dentro de Objetivos y sin
% términos que el capítulo no presentó (revisión del autor, 2026-10-04; tutor, Audio 1).

En este capítulo se introduce la problemática que da origen al presente trabajo, se contextualiza su relevancia
y se describen los objetivos y el enfoque propuesto para abordarla. La Sección~\ref{sec:intro_contexto}
describe la relevancia del cáncer de piel, en el mundo y en la Argentina, las limitaciones que demoran su
diagnóstico y su documentación, y lo que pueden aportar los métodos computacionales, teniendo en cuenta sus
ventajas y desventajas. La Sección~\ref{sec:intro_objetivos} enuncia las preguntas que guían el
trabajo y sus objetivos, y la Sección~\ref{sec:intro_organizacion} describe cómo se organiza el resto del documento.

\section{Contexto y motivación}
\label{sec:intro_contexto}

El cáncer de piel es uno de los tipos de cáncer más frecuentes del mundo. Según la última estimación de
GLOBOCAN, publicada en 2026, durante el año 2024 se diagnosticaron 337\,969 casos nuevos de melanoma y
1\,138\,314 de cáncer de piel no melanoma, que agrupa a aquellos tumores malignos de la piel distintos del melanoma
sin contar el carcinoma basocelular, que la estimación excluye~\cite{sung2026global}. En la Argentina, la misma
estimación registra 1\,602 casos nuevos de melanoma en ese mismo año~\cite{ervik2026cancertoday}.

Sin embargo, no todos los cánceres de piel tienen la misma gravedad. El melanoma representa menos de la cuarta
parte de los diagnósticos, pero causa cerca de la mitad de las muertes por cáncer de piel~\cite{sung2026global}.
En la Argentina, la relación entre muertes y diagnósticos es más desfavorable: las 571 muertes por melanoma
registradas en 2024 equivalen a casi una de cada tres diagnósticos, frente a una de cada seis en el total
mundial~\cite{ervik2026cancertoday}. Además, la enfermedad no se presenta igual en todo el país. Según el
Registro Argentino de Melanoma Cutáneo, en el Noroeste y el Noreste fueron más frecuentes las lesiones de más
de 2~mm de espesor de Breslow, la medida de profundidad del tumor con la que se estima su
pronóstico~\cite{loria2020cutaneous}.

Esa diferencia de espesor importa porque el desenlace del melanoma depende del momento en que se lo
diagnostica. La supervivencia relativa a cinco años es de 100.0\,\% cuando se lo detecta localizado, baja a
76.0\,\% cuando alcanzó los ganglios regionales y a 34.0\,\% cuando hay metástasis a
distancia~\cite{nci2026melanoma}. Por lo tanto, la detección temprana es decisiva y depende de que una lesión
sospechosa sea examinada e interpretada a tiempo por alguien con la capacidad de reconocerla.

Que esa valoración ocurra a tiempo depende, en primer lugar, de la disponibilidad de especialistas. Un
relevamiento en 158 países estimó una media de 2.66 dermatólogos por cada 100\,000 habitantes, con 0.37 en los
países de ingreso bajo y 5.05 en los de ingreso alto~\cite{freeman2026global}. Al mismo tiempo, las lesiones que
requieren una valoración son frecuentes: en la campaña nacional de prevención de 2022 de la Sociedad Argentina
de Dermatología, los dermatólogos encontraron lesiones sospechosas en una de cada cuatro de las 3\,635
personas examinadas en 17 provincias y la Ciudad de Buenos Aires y se registraron 70 lesiones con sospecha de
melanoma~\cite{sad2023campana}. Donde los especialistas no alcanzan, la primera valoración queda en manos de
médicos no especializados en dermatología, lo que dificulta la detección temprana del melanoma.

Quién realiza esa primera valoración condiciona la interpretación de la lesión. La dermatoscopía, el examen de
la lesión con aumento e iluminación controlada, mejora el diagnóstico, pero sobre todo en manos de usuarios con
experiencia~\cite{dinnes2018dermoscopy}. Además, la lectura de una misma imagen no es uniforme, sino que varía
entre observadores según su nivel de formación~\cite{tschandl2019comparison}. En consecuencia, el resultado de
la valoración depende en buena medida de la formación de quien la realiza.

Cada consulta exige, además de diagnosticar, un reporte escrito que registre lo observado y la conclusión.
Redactarlo es una tarea distinta de la interpretación de la lesión, y también consume tiempo: en atención
primaria, los médicos dedican 5.9 horas de una jornada de 11.4 a la historia clínica electrónica, y casi la
mitad de ese tiempo a tareas administrativas y de registro~\cite{arndt2017tethered}. Así, los tiempos de
atención se alargan en dos tareas: la interpretación de la lesión y la redacción de su reporte.

Para asistir en esas dos tareas, se han ido diseñando e implementando diversas herramientas computacionales de
apoyo al diagnóstico. Los primeros métodos analizaban la imagen de la lesión para asignarle una categoría
diagnóstica~\cite{codella2019skin,tschandl2019comparison}, y los más recientes pueden, además, producir un texto
sobre ella~\cite{zhou2023skingpt4}. Con el aprendizaje profundo, las
\glspl{cnn} y, más tarde, los \glspl{vit} aprendieron a clasificar imágenes de lesiones a partir de ejemplos
etiquetados, y en 2017 una \gls{cnn} alcanzó un desempeño comparable al de
dermatólogos~\cite{esteva2017dermatologistlevel}. Con la difusión de los modelos de lenguaje surgieron después
los \glspl{vlm}, que procesan imagen y texto a la vez y pueden tanto asignar una categoría a una imagen como
redactar un texto sobre ella~\cite{radford2021learning,liu2023visual}, con adaptaciones recientes al dominio
médico~\cite{li2023llavamed}. De este modo, las dos tareas, clasificar la lesión y redactar su reporte, pueden
abordarse con un mismo tipo de modelo.

Sin embargo, estos métodos presentan dificultades que todavía limitan su uso clínico. En la clasificación, los
desempeños que informan las revisiones y los estudios recientes se obtienen sobre conjuntos públicos de imágenes
y no en el uso clínico~\cite{magalhaes2024systematic,kreouzi2024learning}, y queda abierto si se sostienen en
atención primaria~\cite{jones2022intelligence} o en otras poblaciones. En particular, el desempeño cae en los
tonos de piel que están poco representados en los datos de entrenamiento~\cite{daneshjou2022disparities,groh2021neural}.
La falta de datos representativos alcanza también a la población argentina, que, con las diferencias regionales
que muestra su registro, recién cuenta desde 2023 con un conjunto público de imágenes
propio~\cite{riccilara2023dataset}.

En la generación de texto aparece una dificultad adicional: un reporte puede ser fluido y aun así afirmar un
diagnóstico equivocado o describir hallazgos que la imagen no muestra~\cite{bai2024hallucination}. Como el
reporte debe reflejar el diagnóstico, las dos tareas están conectadas. Un error de clasificación se traslada al
texto, y a la vez un buen clasificador no garantiza un buen reporte, porque el texto puede describir mal la
lesión aunque el diagnóstico sea correcto.

Estas dificultades dejan dos cuestiones abiertas, que motivan el presente trabajo. La primera es si un
\gls{vlm}, que además de clasificar puede redactar, alcanza en la clasificación de lesiones de piel, y en
particular del melanoma, un desempeño comparable al de los clasificadores entrenados solo con imágenes. La
segunda es si los reportes que genera afirman el diagnóstico correcto y describen la lesión que muestra la
imagen. Para abordarlas, el trabajo estudia en conjunto la clasificación y la generación de reportes, con
recursos computacionales convencionales (\emph{commodity hardware}): cada modelo se entrena y se ejecuta en una
sola GPU, sin requerir infraestructura de cómputo de alto rendimiento, como los nodos con varias GPU de un
centro de datos.


\section{Objetivos}
\label{sec:intro_objetivos}
A partir de la problemática expuesta, el trabajo plantea dos preguntas de investigación que orientan sus objetivos:
\begin{enumerate}[label=\textbf{PI\arabic*.}]
  \item ¿Clasifican los \glspl{vlm} las lesiones de piel, y en particular el melanoma, tan bien como los
  clasificadores entrenados solo con imágenes?
  \item ¿Los reportes que genera un \gls{vlm} contienen el diagnóstico correcto y en qué medida describen la lesión
  que muestra la imagen?
\end{enumerate}

Para responderlas, el objetivo general es desarrollar y evaluar un sistema basado en modelos de aprendizaje
profundo y \glspl{vlm}, adaptados a partir de modelos preentrenados, capaz de asistir en el diagnóstico de lesiones de piel y en la
generación automática de reportes clínicos estructurados a partir de imágenes, con el fin de apoyar el diagnóstico y la documentación de las lesiones. Se analiza el desempeño del sistema en la clasificación, si el reporte contiene el diagnóstico correcto y si puede ejecutarse con recursos computacionales convencionales.

Los objetivos específicos son:
\begin{enumerate}[label=\textbf{OE\arabic*.}]
  \item Diseñar, implementar y evaluar clasificadores de lesiones de piel basados en \glspl{cnn}, \glspl{vit}
  y \glspl{vlm}, generales y especializados en dermatología, para identificar posibles melanomas.
  \item Diseñar, implementar y evaluar un \gls{vlm}, ajustado a partir de un modelo preentrenado, que integre
  información visual y textual para producir un diagnóstico preliminar y el reporte de la lesión.
  \item Evaluar el desempeño del sistema en distintos escenarios, como imágenes de otro conjunto de datos o casos en
  los que las predicciones de los componentes presentan desacuerdo o evidencia insuficiente para afirmar un diagnóstico, analizando su exactitud diagnóstica, la calidad del texto generado y su costo computacional.
\end{enumerate}

La Tabla~\ref{tab:intro_objetivos_preguntas} relaciona cada objetivo específico con la pregunta que contribuye a
responder y con los capítulos donde se lo describe y se lo evalúa.

\begin{table}[H]
  \centering
  \caption{Relación entre los objetivos específicos, las preguntas y los capítulos.}
  \label{tab:intro_objetivos_preguntas}
  \small
  \begin{tabularx}{\textwidth}{@{}>{\raggedright\arraybackslash}p{0.29\textwidth}p{0.10\textwidth}
    >{\raggedright\arraybackslash}X>{\raggedright\arraybackslash}X@{}}
    \toprule
    \textbf{Objetivo específico} & \textbf{Pregunta} & \textbf{Se diseña y describe en} & \textbf{Se evalúa y responde en} \\
    \midrule
    OE1. Clasificadores de lesiones y de melanoma (\gls{cnn}, \gls{vit}, \gls{vlm}) & PI1 &
      Capítulo~\ref{cp:metodos}, Materiales y métodos &
      Capítulo~\ref{cp:res_clasif}, Resultados de clasificación \\
    \filasep
    OE2. \gls{vlm} para generar reportes & PI2 &
      Capítulo~\ref{cp:metodos}, Materiales y métodos &
      Capítulo~\ref{cp:res_gen}, Resultados de generación \\
    \filasep
    OE3. Evaluación en distintos escenarios & PI1 y PI2 &
      Capítulo~\ref{cp:metodos} (protocolo de evaluación) &
      Capítulos~\ref{cp:res_clasif} y~\ref{cp:res_gen} (resultados) y Capítulo~\ref{cp:discusion} (discusión y
      limitaciones) \\
    \bottomrule
  \end{tabularx}
\end{table}

\section{Organización del documento}
\label{sec:intro_organizacion}
Después de esta introducción, el documento se organiza en seis capítulos:

\begin{description}
  \item[Capítulo~\ref{cp:marco}, Marco teórico.] Este capítulo presenta los conceptos necesarios para leer el
  resto del documento y revisa los trabajos más cercanos.
  \item[Capítulo~\ref{cp:metodos}, Materiales y métodos.] Este capítulo presenta los conjuntos de datos, los
  modelos, el \emph{pipeline} que encadena la clasificación y la generación del reporte, el protocolo de
  evaluación y el diseño de los experimentos.
  \item[Capítulo~\ref{cp:res_clasif}, Resultados de clasificación.] Este capítulo presenta la comparación de los
  clasificadores supervisados y de los \glspl{vlm}, la detección de melanoma y la evaluación del clasificador elegido
  con imágenes de otro conjunto de datos y con imágenes degradadas, junto con su calibración y su costo, y la comparación con el estado del arte.
  \item[Capítulo~\ref{cp:res_gen}, Resultados de generación.] Este capítulo presenta la evaluación de los reportes
  a lo largo del \emph{pipeline}: el generador sin diagnóstico, la cascada con el clasificador, el ruteo por
  concordancia, la comparación entre generadores, sus costos computacionales y el contraste con el estado del arte.
  \item[Capítulo~\ref{cp:discusion}, Discusión.] Este capítulo presenta la interpretación conjunta de los
  resultados y las limitaciones del trabajo.
  \item[Capítulo~\ref{cp:conclusiones}, Conclusiones.] Este capítulo presenta la respuesta a las preguntas
  planteadas, los hallazgos principales, el
  balance de los objetivos, los aportes del trabajo y el trabajo futuro.
\end{description}

Por último, dos apéndices reúnen el material que, por ser muy detallado, no se incluye en el cuerpo de los
capítulos: los \emph{prompts} y las plantillas que usan los modelos, y la información necesaria para reproducir los
experimentos.

